\section{Preliminaries}
\label{sec:prelim}

\subsection{Notation}

We consider a transformer-based large language model (LLM) with $L$ layers and $H$ attention heads per layer. We denote the set of all attention heads as
$
\mathcal{H} = \{h_{\ell,a} \mid \ell = 1, \ldots, L,\; a = 1, \ldots, H\},
$
where $h_{\ell,a}$ represents the $a$-th head in layer $\ell$. We use $\mathcal{G}=(\mathcal{V},\mathcal{E})$ to denote a graph, where $\mathcal{V}$ is the node set and $\mathcal{E}$ is the edge set. In our functional graphs, nodes correspond to attention heads, i.e., $\mathcal{V}=\mathcal{H}$, and edges are constructed from activation-based similarity between heads.

\subsection{Small-World Organization of Brain Networks}

Small-world organization describes a network topology characterized by high local clustering and efficient global communication \cite{barbey2018network, rosenberg2020functional}. In neuroscience, functional brain networks are often built over anatomical regions of interest (ROIs), with nodes as brain regions and edges as functional connectivity \cite{langer2012functional}. Such organization has been linked to human intelligence: higher Raven's Advanced Progressive Matrices scores \cite{bors2003effect, raven1962advanced} are associated with stronger small-world structure, reflected by higher clustering coefficients and shorter communication path lengths \cite{langer2012functional}.

More generally, given a graph $\mathcal{G}=(\mathcal{V},\mathcal{E})$, let $C(\mathcal{G})$ denote its average clustering coefficient \cite{watts1998collective}, defined as
\begin{equation}
C(\mathcal{G}) =
\frac{1}{|\mathcal{V}|}
\sum_{i\in\mathcal{V}} C_i,
\qquad
C_i =
\frac{2e_i}{d_i(d_i-1)},
\end{equation}
where $d_i$ is the degree of node $i$, $e_i$ is the number of edges among the neighbors of node $i$, and we set $C_i=0$ when $d_i<2$. We let $E(\mathcal{G})$ denote its global efficiency, which captures efficient communication through short paths \cite{latora2001efficient}:
\begin{equation}
E(\mathcal{G}) =
\frac{1}{|\mathcal{V}|(|\mathcal{V}|-1)}
\sum_{\substack{i,j\in\mathcal{V}\\ i\neq j}}
\frac{1}{\operatorname{dist}(i,j)},
\end{equation}
where $\operatorname{dist}(i,j)$ is the shortest-path distance between nodes $i$ and $j$, and disconnected node pairs contribute zero efficiency. To assess small-world organization, $\mathcal{G}$ is compared against degree-preserving randomized graphs generated by edge swaps \cite{hosseini2013influence}. Let $C_{\mathrm{rand}}$ and $E_{\mathrm{rand}}$ denote the average clustering coefficient and global efficiency of these randomized graphs. The small-world index of $\mathcal{G}$ is then defined as
\begin{equation}
\mathrm{SWI}(\mathcal{G}) =
\frac{C(\mathcal{G})}{C_{\mathrm{rand}}}
\cdot
\frac{E(\mathcal{G})}{E_{\mathrm{rand}}}.
\label{eq:swi}
\end{equation}
A larger SWI indicates stronger local clustering and more efficient global communication relative to degree-matched random graphs.


\subsection{SVCCA for Representation Similarity}
\label{sec:svcca_prelim}

We use Singular Vector Canonical Correlation Analysis (SVCCA)~\cite{raghu2017svcca} to measure similarity between vector-valued representations. Given two mean-centered representation matrices
$\mathbf{X}\in\mathbb{R}^{m\times D_x}$ and
$\mathbf{Y}\in\mathbb{R}^{m\times D_y}$ over the same $m$ examples, SVCCA first applies singular value decomposition:
\begin{equation}
\mathbf{X} = \mathbf{U}_X \mathbf{S}_X \mathbf{V}_X^\top,
\qquad
\mathbf{Y} = \mathbf{U}_Y \mathbf{S}_Y \mathbf{V}_Y^\top .
\end{equation}

We retain the top singular directions explaining a high proportion of variance, e.g., $99\%$. Let $p_x$ and $p_y$ be the numbers of retained directions for $\mathbf{X}$ and $\mathbf{Y}$, respectively. The representations are then projected as

\begin{equation}
\mathbf{X}' = \mathbf{X}\mathbf{V}_X^{(p_x)}, 
\qquad
\mathbf{Y}' = \mathbf{Y}\mathbf{V}_Y^{(p_y)} ,
\end{equation}

where $\mathbf{X}'\in\mathbb{R}^{m\times p_x}$ and
$\mathbf{Y}'\in\mathbb{R}^{m\times p_y}$. Here, $\mathbf{V}_X^{(p_x)}$ and $\mathbf{V}_Y^{(p_y)}$
denote the matrices formed by the first $p_x$ and $p_y$ columns of
$\mathbf{V}_X$ and $\mathbf{V}_Y$, respectively.

Canonical Correlation Analysis (CCA)~\cite{thompson2000canonical} then finds projection vectors $\mathbf{w}_x$ and $\mathbf{w}_y$ that maximize the correlation between the projected representations:
$
\rho =
\max_{\mathbf{w}_x, \mathbf{w}_y}
\frac{
\mathbf{w}_x^\top \boldsymbol{\Sigma}_{X'Y'} \mathbf{w}_y
}{
\sqrt{
(\mathbf{w}_x^\top \boldsymbol{\Sigma}_{X'X'} \mathbf{w}_x)
(\mathbf{w}_y^\top \boldsymbol{\Sigma}_{Y'Y'} \mathbf{w}_y)
}
}.
$
Here, $\boldsymbol{\Sigma}_{X'X'}$ and $\boldsymbol{\Sigma}_{Y'Y'}$ denote the covariance matrices of
$\mathbf{X}'$ and $\mathbf{Y}'$, and $\boldsymbol{\Sigma}_{X'Y'}$ denotes their cross-covariance matrix. This produces a set of canonical correlations
$\{\rho_q\}_{q=1}^{r}$, where $r=\min(p_x,p_y)$.
We define the SVCCA similarity as their mean:
\begin{equation}
\mathrm{SVCCA}(\mathbf{X},\mathbf{Y})
=
\frac{1}{r}\sum_{q=1}^{r}\rho_q .
\end{equation}
Higher SVCCA values indicate stronger alignment between representational subspaces.